To reduce wildfire risk during windy and dry conditions, electrical utility companies have started to preemptively shut off power. While intended to protect public safety, these public safety power shutoffs (PSPS) may themselves affect health—particularly for people with respiratory conditions or dependence on electricity-powered medical equipment. Using emergency department and hospital records across California from 2013–2019, we found that severe PSPS events increased the odds of respiratory healthcare visits, with stronger effects among people with chronic obstructive pulmonary disease. When PSPS events co-occurred with wildfire smoke exposure, respiratory risks were even higher, suggesting interaction that may compound harm. As climate change drives more frequent wildfires, utilities and policymakers must not only weigh the health costs of preemptive PSPS events, but also invest in preparedness strategies—such as public health messaging, accessible cooling and charging centers, and support for medically vulnerable residents—to reduce harm when shutoffs do occur.